\documentclass[conference]{IEEEtran}
\IEEEoverridecommandlockouts
\usepackage{cite}
\usepackage{amsmath}
\usepackage{booktabs}
\usepackage{url}
\usepackage[hidelinks]{hyperref}

% Reproducible output, and the reason this is three primitives rather than an
% environment variable. SOURCE_DATE_EPOCH pins the creation date, but pdfTeX
% also derives the trailer /ID from the build path, so the same source in a
% different directory produced a PDF that differed in exactly those 32 bytes --
% identical content, identical length, identical xref offsets, different ID.
% That is what made the committed paper fail its own byte-comparison test in a
% fresh clone while passing in the directory it was built in.
\ifdefined\pdfinfoomitdate \pdfinfoomitdate=1 \fi
\ifdefined\pdftrailerid \pdftrailerid{} \fi
\ifdefined\pdfsuppressptexinfo \pdfsuppressptexinfo=-1 \fi

\begin{document}

\title{Verifying Retention Analytics Against a Known Hazard Process:\\
An Engineering Case Study}

\author{\IEEEauthorblockN{Saurabh Mastud}
\IEEEauthorblockA{Independent project\\
\texttt{github.com/SaurabhMastud/churnlens}}}

\maketitle

\begin{abstract}
Cohort retention tables, Kaplan--Meier curves, churn-driver models and
lifetime-value estimates are routine analytics deliverables, and they share a
structural weakness: computed on observational data, a correct analysis and a
plausible-looking incorrect one produce the same artifact. This paper describes
ChurnLens, a seven-day engineering exercise that inverts the usual setup.
Subscriber lifecycles are generated from an explicit discrete-time hazard model
whose coefficients are fixed in configuration and stored beside the data in
tables that analysis code is structurally prevented from reading. Every
analysis can then be asked whether it recovers the process that produced its
input. We report the architecture, the two-level enforcement of that
separation, and a test taxonomy that distinguishes exact bookkeeping assertions
from statistical comparisons with stated tolerances. The evaluation is an
honest engineering case study rather than a controlled experiment: we enumerate
the seven defects that reached a committed state, their root causes, and what
caught each. Three are of general interest --- a generated target leak that
would have made a leakage artifact look like a near-perfect model, an optimiser
default that stopped short of the maximum likelihood estimate and presented as
a coefficient-recovery failure, and a non-deterministic quantile assignment
that made an analysis irreproducible between runs on the same data. We close
with threats to validity, including the fact that the generator hands the
driver model its own functional form.
\end{abstract}

\begin{IEEEkeywords}
survival analysis, data quality, software testing, churn prediction, measurement
error, reproducibility
\end{IEEEkeywords}

\section{Introduction}

A retention analysis is a claim about an unobserved process. Published examples
are almost always built on observational data whose true churn mechanism nobody
knows, so the analysis can be sanity-checked but not verified: an incorrectly
handled censoring rule, a quantile defined on post-outcome information, or a
leaked target produce output that is indistinguishable, at the artifact level,
from output produced correctly. The failure mode that matters most is the one
that looks like success.

ChurnLens addresses that gap by construction. Subscriber lifecycles are drawn
from a discrete-time hazard
\begin{equation}
\mathrm{logit}\,h(t) = \beta_0 + \beta_{\mathrm{plan}} + \beta_\tau \log(1+t)
+ \beta_e e_i
\label{eq:hazard}
\end{equation}
with coefficients fixed in configuration, where $t$ is tenure in months and
$e_i \in [0,1]$ is a latent per-subscriber engagement level. The coefficients
and the latent levels are written into the store alongside the generated data,
in tables the analysis modules cannot read. Tests may read them. Every estimate
therefore has a ground-truth counterpart.

This paper claims the following, and nothing beyond it:

\begin{enumerate}
\item A one-way ground-truth channel --- truth readable by tests, structurally
unreachable by analyses --- is implementable in a small project and catches
analytical errors that no amount of internal consistency checking would.
\item Separating that channel \emph{structurally} (truth is not in the queried
tables) dominates separating it \emph{by convention or by test}, and we show
the specific hole the test-only version left.
\item Retention, survival, driver-model and value estimates computed from
observable columns alone land within stated tolerances of the planted process,
and we report the one quantity that does not: the engagement coefficient, which
attenuates to 58\,\% of its planted value for reasons that are predictable and
measurable rather than a defect.
\item Seven defects reached a committed state over seven days. We give each
one's root cause and the mechanism that caught it, including three that would
have invalidated the analysis while leaving the test suite green.
\end{enumerate}

We do not claim predictive performance, generality beyond synthetic data, or
that the approach is novel in a statistical sense; the components are standard
and cited below. The contribution is an engineering account of making
analytical correctness testable.

\section{Background and Related Work}

The estimators are classical. Kaplan and Meier's product-limit
estimator~\cite{kaplan1958} handles right-censored observation; Greenwood's
formula~\cite{greenwood1926} supplies its variance, which we use on the
complementary log-log scale so that interval bounds remain inside $[0,1]$.
Group comparison uses the log-rank statistic~\cite{mantel1966}. Because the
data is monthly by construction, the driver model is a discrete-time hazard
model in the sense of Allison~\cite{allison1982}, which makes logistic
regression on subscriber-months the correctly specified estimator for
(\ref{eq:hazard}) rather than an approximation of a continuous-time one.

Three literatures inform the specific errors we instrument. Measurement error
in a regressor attenuates its estimated coefficient toward zero by a reliability
ratio~\cite{fuller1987}; we observe this directly, since the generator drives
churn from a latent level while the analysis sees only noisy monthly event
counts. Defining a group membership using information generated after the start
of observation is immortal-time bias~\cite{suissa2008}; the natural phrasing of
``segment subscribers by their average activity'' commits exactly that error.
Leakage as a software-engineering phenomenon --- a feature that encodes the
target, or an evaluation split that shares entities across folds --- is
characterised by Kaufman \emph{et al.}~\cite{kaufman2012}.

On the reporting side, we follow Gneiting and Raftery~\cite{gneiting2007} in
using a proper scoring rule rather than accuracy, since the monthly churn rate
is below 2\,\% and a constant prediction is 98\,\% accurate; calibration is
assessed separately from discrimination following standard
practice~\cite{harrell2015}. For lifetime value, Fader and
Hardie~\cite{fader2010} show that applying an aggregate churn rate to a
heterogeneous base misstates customer value; our hazard both declines with
tenure and varies across subscribers, so the familiar
$\mathrm{ARPU}/\mathrm{churn}$ shortcut is wrong here in a way we can quantify
exactly. Implementation rests on scikit-learn~\cite{pedregosa2011} and
DuckDB~\cite{raasveldt2019}.

\section{System Architecture}

\begin{table}[t]
\caption{Components. All run; none is aspirational.}
\label{tab:components}
\centering
\begin{tabular}{@{}lll@{}}
\toprule
Component & Tool & Role \\
\midrule
Hazard model & Python & Coefficients; ground truth \\
Generator & Python, NumPy & Lifecycles with right-censoring \\
Store & DuckDB & 4 analysis tables, 2 quarantined \\
Sanity report & Python, SQL & Churn by plan, hazard by tenure \\
Cohorts & Python, SQL & Retention matrix, pooled curves \\
Survival & Python, SQL & Kaplan--Meier, log-rank \\
Drivers & scikit-learn & Coefficients, calibration, scores \\
Segments & Python, SQL & Activity-segment survival \\
Value & Python & Restricted expected value \\
Dashboard & Streamlit & Five tabs over the above \\
\bottomrule
\end{tabular}
\end{table}

Table~\ref{tab:components} lists the components. The data flow is linear: the
hazard configuration feeds the generator, which writes a DuckDB store; every
analysis reads that store; the dashboard and the reports read the analyses.

Three decisions are load-bearing.

\textbf{The store is split by the contract.} Four tables are analysis-facing
(\texttt{subscribers}, \texttt{subscription\_months},
\texttt{engagement\_events}, \texttt{run\_meta}); two are quarantined
(\texttt{hazard\_truth}, holding the coefficients, and
\texttt{subscriber\_truth}, holding each subscriber's latent engagement level).
Analyses may read the first group only.

\textbf{Enforcement is structural first and by test second.} On day~1 the
latent engagement level sat on an analysis-facing table and only a per-module
test guarded it. That was not a theoretical weakness: substituting the latent
level for observed activity passed every other test in the file and produced a
cleaner result, because it was reading the answer. Day~2 moved the column to
\texttt{subscriber\_truth}, after which an analysis cannot read what is not in
the table it queries. The test-based guard remains as a second line, and now
discovers analysis modules by convention --- anything under \texttt{src/} or
\texttt{dashboard/} that is not on the generator side --- rather than being
opted in one module at a time.

\textbf{The observation window is recorded, not inferred.} Cohort exposure
determines which cells of a retention table are observable, and the natural
derivation (\texttt{max(signup + tenure) + 1}) is exact only if some subscriber
is observed to the final month. A one-subscriber 24-month store whose member
signed up in month~15 and churned at tenure~4 reports a 20-month window,
shrinking that cohort's exposure with nothing to indicate error. The window is
therefore carried in a provenance row, and a missing row raises rather than
defaulting.

\section{Implementation}

The repository holds 2{,}765 lines of application and dashboard Python and
3{,}917 lines of tests, a ratio of about 1.4 test lines per source line, built
over seven daily sessions; its history held thirty commits when this paper was
rendered. Analyses are plain functions
taking a DuckDB connection and returning lists of dictionaries; each module
also renders a text report and exposes a command-line entry point, so every
figure quoted in this paper is reproducible from a clean clone with
\texttt{python -m src.<module>}.

Four points of implementation substance:

\emph{Exposure handling.} A retention cell a cohort has not had time to reach
is omitted, never zeroed. Zero-filling is the standard way this goes wrong, and
it is not visually obvious: on the committed store it drags pooled retention at
tenure~12 from 82.5\,\% to 40.5\,\%, which reads as a finding and is an
artifact. The exposure rule exists once, in one parameterised SQL query serving
the matrix, the pooled curve and the segmented curves.

\emph{Censoring.} The survival risk table derives who leaves the risk set from
each subscriber's last observed month rather than from the generator's
censoring flag, so the two can be reconciled as independent sources; they agree
exactly. The bookkeeping is pinned by the identity $n_{t+1} = n_t - d_t - c_t$.

\emph{Model fitting.} The driver model is unpenalised, because the comparison of
interest is against planted coefficients and scikit-learn's default $L_2$
penalty would bias the estimates by an amount nobody chose. Standard errors come
from the inverse Fisher information, computed over the design matrix the fit
already built, which avoids a second modelling dependency.

\emph{Restricted value.} Expected value is computed over an explicit horizon
and a horizon longer than the observed curve is refused. Survival never reaches
zero inside the window --- it bottoms out at 75.9\,\% --- so an unrestricted
``lifetime'' figure would be an extrapolation of the tail presented as a
measurement.

\section{Verification Strategy}

The suite contains 188 test functions in fourteen modules, 249 cases after
parametrisation. It is organised around a distinction we found useful: tests of
bookkeeping are exact, tests of estimates are statistical, and mixing them
produces a suite that is loose where it could be strict.

\begin{enumerate}
\item \textbf{Exact invariants.} The risk-set identity; a cohort's retention
matrix reconciled cell by cell against the subscriber-month rows; segmented
risk tables summing to the pooled one; censoring counts matching the
generator's flags; the design matrix retaining every subscriber-month.
\item \textbf{Estimate versus planted truth}, with a stated tolerance and a
stated population. Cohort retention and the Kaplan--Meier curve agree with the
planted survival function within four standard errors at every tenure checked;
the tenure-0 hazard comes out at 3.833\,\% observed against 3.920\,\%
predicted, a gap of 0.34 binomial standard errors; restricted value agrees
within 2\,\%.
\item \textbf{Calibration of the inferential tooling itself.} The log-rank test
is run against a dataset rebuilt with every plan effect set to zero, so the
null is true by construction. Across ten seeds it rejected once, consistent
with its nominal 5\,\% size. The assertion is on the rejection rate, not on a
single draw --- demanding zero rejections from a correctly sized test would
assert that the test is broken.
\item \textbf{Mutation checks.} Each guard is verified to fail when the
behaviour it protects is reverted. Tripling every predicted hazard leaves the
area under the ROC curve identical to twelve decimal places and moves the worst
calibration decile past ten standard errors, which is what establishes that the
calibration test is doing work and that discrimination alone would not notice.
\item \textbf{Execution, not just import.} The dashboard is run headless
through Streamlit's own script-run harness. We first used a health-endpoint
probe, which returns \texttt{ok} in about three seconds and proves only that
the web server started: Streamlit executes the page script per session, so with
no client attached a page that would crash on its first line still answers 200.
\end{enumerate}

Two negative results belong here because they bound what the suite can claim.
First, predictive scoring is a weak instrument for specification error: the
correctly specified model beats a linear-tenure misspecification on held-out
log loss in all five seeds tried, by 0.4\,\%. No practitioner would notice that
margin, while the misspecified coefficient means something materially different.
Second, the textbook warning about splitting panel data by row rather than by
entity is directionally right but small here --- 0.0025 in mean area under the
curve across five seeds, against a between-seed standard deviation of 0.014 ---
because a four-parameter model has nothing to memorise. The split is still by
subscriber, since it is free and the harm scales with model capacity, but the
test asserts the measured claim rather than the received one.

\section{Evaluation}

This is an engineering case study, not a controlled experiment: there is no
baseline implementation to compare against and no human subjects. The evidence
we can offer honestly is the defect record. Table~\ref{tab:defects} lists every
defect that reached a committed state, with the mechanism that caught it.

\begin{table*}[t]
\caption{Defects that reached a committed state, with root cause and detector.}
\label{tab:defects}
\centering
\begin{tabular}{@{}llll@{}}
\toprule
\# & Defect & Root cause & Caught by \\
\midrule
1 & Churn months had no engagement row, making ``no activity'' a perfect
target separator & Two copies of a span expression that drifted apart & Profiling
the suite, then inspecting row counts \\
2 & Planted coefficients were mutable despite a frozen dataclass & \texttt{frozen}
blocks rebinding, not mutation of a contained dict & Mutation test written
against the guard \\
3 & Dependency pins described an environment never installed & File written from
memory, never installed from & Comparing pins against the live interpreter \\
4 & Latent engagement readable from an analysis-facing table & Contract enforced
per module by test rather than structurally & Mutation-testing a new analysis
cut \\
5 & Optimiser stopped short of the maximum likelihood estimate & Library default
tolerance of $10^{-4}$ & Score identity $\sum(y_i - \hat p_i) = 0$ \\
6 & Dashboard crashed when a slider exceeded the observed window & Slider bounds
hardcoded; empty result pivoted & Running the app headless \\
7 & Contract guard rejected honest prose containing a column name & Guard matched
every string literal, not just queries & Writing a user-facing caption \\
\bottomrule
\end{tabular}
\end{table*}

Three warrant comment.

\textbf{Defect 1} is the one the project exists to catch. The generator gave a
churned subscriber a subscription-month row for the month they left but no
engagement row for it, so all 162 churn rows lacked an activity record and all
10{,}228 others had one. An inner join would have dropped every churn event,
which fails loudly; a left join with zero-filling would have made ``no
activity'' predict churn exactly and reported a leakage artifact as a
near-perfect model. Both tables now derive their span from one shared function.

\textbf{Defect 5} is the most instructive. At scikit-learn's default tolerance,
L-BFGS left a residual of 0.98 events on 43{,}467 rows, and returned a plan
coefficient of $-0.425$ where the converged fit gives $-0.4516$ against a
planted $-0.45$. The symptom was a coefficient-recovery failure --- precisely
the signal this project is built to treat as a bug in the analysis --- and the
cause was the optimiser quitting early. The detector is worth noting because it
needs no ground truth: the score equation for the intercept requires fitted
probabilities to sum to the observed event count, which is checkable on real
data.

\textbf{Attenuation is not in the table}, because it is not a defect. Fitted on
the latent regressor the model recovers every coefficient within 1.5 standard
errors; fitted on observable activity the engagement coefficient comes back at
$-1.383$ against a planted $-2.40$, which is 8.3 standard errors away and
therefore bias rather than noise. The reliability ratio
$\sigma^2_e/(\sigma^2_e + \overline{e}/12) = 0.0489/(0.0489+0.0427) = 0.53$
predicts the 0.58 observed; the residual gap is attributable to the nonlinear
link and to the proxy being a single month rather than a subscriber average.
Reporting $-1.4$ as ``the engagement effect'' would understate the strongest
lever in the modelled business by nearly half, and nothing in the fit's own
diagnostics would say so.

Two further defects were caught before any commit and are recorded for
completeness: a quantile assignment made non-deterministic by ties at the bin
boundary (the same store yielded 66.1\,\% and then 65.9\,\% for one cell), and
an inferred observation window that silently shrank a cohort's exposure.

\section{Discussion and Threats to Validity}

\subsection{Lessons}

A contract maintained by a test that someone must remember to extend is one
module away from being maintained by nothing; discovery by convention is
cheaper than vigilance. A guard that forces the product to get worse in order
to stay green --- here, a caption that had to avoid the clearest word for what
it described --- has a bug in the guard, and narrowing its scope with a mutation
test is better than rewording the output. Finally, several of the most useful
results in this project were negative: the measured size of a textbook bias,
the measured weakness of held-out scoring as a specification check.

\subsection{Construct validity}

The quantity we verify is ``does this estimator recover the coefficients of the
process that generated its input''. That is narrower than ``is this analysis
useful''. An analysis can recover the planted process and still answer the wrong
business question, and nothing here addresses that.

\subsection{Internal validity}

Most tolerances are stated in standard errors, but some are absolute (2\,\% for
restricted value, 0.01 for the tenure-0 hazard) and were chosen to be loose
enough to survive a reseed. Where a bound was sensitive we measured it across
seeds and recorded the range rather than quoting the first value found; the
look-ahead inflation figure, for instance, ranges 1.13--1.54 across five store
configurations, and the single number initially observed was the largest.
Statistical tests on the same store are not independent of one another.

\subsection{External validity}

The data is synthetic. It has no missing values, no duplicate records, no plan
changes mid-life, no reactivation, no seasonality and no concept drift --- all of
which a real retention dataset has and all of which are the usual reason an
analysis is hard. More pointedly, the generator hands the driver model its own
functional form: the hazard is logit-linear in $\log(1+t)$ and the model fits
exactly that, so coefficient recovery is a test of the implementation rather
than of model selection. The one misspecification we examine is a linear tenure
term. Results here should be read as evidence that the implementation is
correct, not that the approach transfers unchanged to observational data.

\subsection{Unexercised components}

Every component in Table~\ref{tab:components} runs, and the dashboard is
executed in the test suite rather than merely imported. Three gaps remain. The
dashboard has never been used by anyone other than its author, so its captions
are untested as communication. The generator's behaviour at extreme scale is
unmeasured: the largest store built is 6{,}000 subscribers over 36 months. And
the log-rank implementation uses the conservative $\sum(O-E)^2/E$ form rather
than the full variance expression, so a non-significant result from it should be
read as ``not shown'' rather than ``no difference'' --- which the calibration run
confirmed, finding zero rejections in twenty null datasets at 2{,}000
subscribers.

\section{Conclusion and Future Work}

Generating data from a known hazard process and quarantining that process from
the analysis made analytical correctness testable in a way that internal
consistency checks cannot be, and the defect record suggests the cost was
repaid: four of the seven committed defects were errors of inference that a
conventional test suite would not have been positioned to notice. The two
enforcement levels are not equivalent, and the structural one should come first.

Three extensions follow directly. Injecting known pathologies --- missing
engagement rows, duplicated subscribers, a mid-life plan change --- would let
the suite assert that an analysis degrades detectably rather than silently.
Deliberately misspecifying the generator relative to the model would test model
selection rather than only implementation. And a correction for the measured
attenuation, such as fitting on a subscriber-level activity average with its
reliability estimated from the within-subscriber variance, would turn the day-4
negative result into a method with a ground-truth check already in place.

\begin{thebibliography}{00}

\bibitem{kaplan1958} E. L. Kaplan and P. Meier, ``Nonparametric estimation from
incomplete observations,'' \emph{J. Amer. Statist. Assoc.}, vol. 53, no. 282,
pp. 457--481, 1958.

\bibitem{greenwood1926} M. Greenwood, ``The natural duration of cancer,''
\emph{Reports on Public Health and Medical Subjects}, no. 33, pp. 1--26, London:
HMSO, 1926.

\bibitem{mantel1966} N. Mantel, ``Evaluation of survival data and two new rank
order statistics arising in its consideration,'' \emph{Cancer Chemotherapy
Reports}, vol. 50, no. 3, pp. 163--170, 1966.

\bibitem{allison1982} P. D. Allison, ``Discrete-time methods for the analysis of
event histories,'' \emph{Sociological Methodology}, vol. 13, pp. 61--98, 1982.

\bibitem{fuller1987} W. A. Fuller, \emph{Measurement Error Models}. New York:
Wiley, 1987.

\bibitem{suissa2008} S. Suissa, ``Immortal time bias in pharmacoepidemiology,''
\emph{Amer. J. Epidemiology}, vol. 167, no. 4, pp. 492--499, 2008.

\bibitem{kaufman2012} S. Kaufman, S. Rosset, C. Perlich, and O. Stitelman,
``Leakage in data mining: Formulation, detection, and avoidance,'' \emph{ACM
Trans. Knowl. Discov. Data}, vol. 6, no. 4, article 15, 2012.

\bibitem{gneiting2007} T. Gneiting and A. E. Raftery, ``Strictly proper scoring
rules, prediction, and estimation,'' \emph{J. Amer. Statist. Assoc.}, vol. 102,
no. 477, pp. 359--378, 2007.

\bibitem{harrell2015} F. E. Harrell, \emph{Regression Modeling Strategies},
2nd ed. Cham: Springer, 2015.

\bibitem{fader2010} P. S. Fader and B. G. S. Hardie, ``Customer-base valuation
in a contractual setting: The perils of ignoring heterogeneity,'' \emph{Marketing
Science}, vol. 29, no. 1, pp. 85--93, 2010.

\bibitem{pedregosa2011} F. Pedregosa \emph{et al.}, ``Scikit-learn: Machine
learning in Python,'' \emph{J. Mach. Learn. Res.}, vol. 12, pp. 2825--2830, 2011.

\bibitem{raasveldt2019} M. Raasveldt and H. M\"uhleisen, ``DuckDB: An embeddable
analytical database,'' in \emph{Proc. 2019 Int. Conf. Management of Data
(SIGMOD)}, pp. 1981--1984, 2019.

\end{thebibliography}

\end{document}
